% Отчёт по НИРС, осень 2026
% Класс BMSTU-IU8 v1.9 (CatInCosmicSpace, MIT) лежит в корне проекта.
% Сборка: pdfLaTeX + biber (см. README.md).
\documentclass[research]{BMSTU-IU8}

\usepackage{amssymb} % \mathbb и прочие символы (класс их не подключает)

\addbibresource{main.bib}

% ---------- Титульный лист (ЗАГЛУШКИ — заменить на реальные данные) ----------
\student{А. А. Смирнова}
\studentFullName{Смирнова Анна Андреевна}
\group{ИУ8-74}
\supervisor{Б. Б. Кузнецов}
\supervisorWithDegree{доцент, к.т.н. Кузнецов Б. Б.}
\speciality{10.05.01 <<Компьютерная безопасность>>}
\specialization{10.05.01\_01 <<Математические методы защиты информации>>}
% Для НИРС тема заполняется с \hfill, чтобы подчёркивание шло до конца строки
\theme{Проектирование защищённой системы машинного обучения
  с~применением технологий многосторонних конфиденциальных вычислений \hfill}

% ---------- Термины и сокращения ----------
\newglossaryentry{mpc-term}{
    name={Многосторонние конфиденциальные вычисления},
    description={совместное вычисление заданной функции, при котором сведения, доступные нарушителю, ограничены его входами, разрешёнными выходами и публичными параметрами в принятой модели безопасности}
}

\newglossaryentry{secret-sharing}{
    name={Разделение секрета},
    description={представление значения набором долей, позволяющее восстановить значение при объединении долей уполномоченного множества участников}
}

\newglossaryentry{replicated}{
    name={Реплицированное разделение},
    description={способ размещения долей секрета, при котором одну компоненту хранят несколько сторон. В рассматриваемой схеме каждый из трёх серверов хранит две из трёх аддитивных компонент}
}

\newglossaryentry{semi-honest}{
    name={Пассивный (получестный) нарушитель},
    description={участник, который следует протоколу и использует доступные ему доли, случайные значения и сообщения для получения сведений о чужих входах}
}

\newglossaryentry{static}{
    name={Статический нарушитель},
    description={нарушитель, множество скомпрометированных сторон которого фиксируется до начала исполнения протокола}
}

\newglossaryentry{threshold}{
    name={Порог компрометации},
    description={максимальное число совместно наблюдаемых нарушителем серверов, при котором сохраняется рассматриваемая гарантия. В принятой модели пассивного нарушителя порог равен одному серверу из трёх}
}

\newglossaryentry{decision-tree}{
    name={Решающее дерево},
    description={модель классификации, в которой внутренние узлы задают условия выбора ветви, а листья содержат метки классов}
}

\newglossaryentry{private-inference}{
    name={Конфиденциальный вывод},
    description={вычисление прогноза по заданной модели с защитой закрытых входов и промежуточных значений и с выдачей результата назначенному получателю}
}

\newglossaryentry{fixed-point}{
    name={Кодирование с фиксированной точкой},
    description={представление числа целым кодом, полученным умножением на общий масштаб и округлением. Для признаков и порогов в работе применяется масштаб $2^f$}
}

\newglossaryentry{reference}{
    name={Целочисленный эталон},
    description={открытый обход дерева по закодированным признакам и порогам с тем же правилом сравнения, которое используется при конфиденциальном выводе}
}

\newacronym{mkv}{МКВ}{многосторонние конфиденциальные вычисления}
\newacronym{aby}{ABY3}{трёхсторонний протокол с арифметическим и двоичным представлениями и схемами Яо (Arithmetic, Binary, Yao)}
\newacronym{tfe}{TF Encrypted}{платформа конфиденциальных вычислений на основе TensorFlow}

\newglossaryentry{sym-tree}{
    type=\acronymtype, name={$T$}, sort={tree},
    description={решающее дерево}
}

\newglossaryentry{sym-input}{
    type=\acronymtype, name={$x$}, sort={input},
    description={вектор признаков объекта}
}

\newglossaryentry{sym-depth}{
    type=\acronymtype, name={$d$}, sort={depth},
    description={глубина дерева в рёбрах}
}

\newglossaryentry{sym-nodes}{
    type=\acronymtype, name={$m$}, sort={nodes},
    description={число внутренних узлов дерева}
}

\newglossaryentry{sym-batch}{
    type=\acronymtype, name={$B$}, sort={batch},
    description={число объектов в пакете}
}

\newglossaryentry{sym-precision}{
    type=\acronymtype, name={$f$}, sort={precision},
    description={число дробных разрядов}
}

\newglossaryentry{sym-impl}{
    type=\acronymtype, name={$e_{\mathrm{impl}}$}, sort={impl},
    description={доля расхождений между конфиденциальным выводом и целочисленным эталоном}
}

\newglossaryentry{sym-quant}{
    type=\acronymtype, name={$e_{\mathrm{quant}}$}, sort={quant},
    description={доля расхождений между целочисленным эталоном и открытым выводом по исходному дереву}
}


\begin{document}
    \maketitle

    \abstract

МНОГОСТОРОННИЕ КОНФИДЕНЦИАЛЬНЫЕ ВЫЧИСЛЕНИЯ, МАШИННОЕ ОБУЧЕНИЕ,
РАЗДЕЛЕНИЕ СЕКРЕТА, РЕШАЮЩЕЕ ДЕРЕВО, ABY3, TF ENCRYPTED, МОДЕЛЬ НАРУШИТЕЛЯ

% ЧЕРНОВИК: заполнить после написания основных глав.
Объектом исследования является система машинного обучения, обрабатывающая
конфиденциальные данные нескольких владельцев.

Цель работы~--- спроектировать защищённую систему машинного обучения на основе
протоколов многосторонних конфиденциальных вычислений и проверить
работоспособность выбранной платформы.

В ходе работы выполнено сопоставление угроз с возможностями МКВ, построена
модель нарушителя, разработаны архитектура и требования, обоснован выбор
протокола ABY3 и платформы TF Encrypted, разработаны методика экспериментов и программа
пробного запуска с аналитически определёнными контрольными результатами.


    \tableofcontents
    \termsanddefenitions
    \listofabbreviations

    \introduction

% ЧЕРНОВИК
Совместное обучение моделей на данных нескольких организаций ограничено
требованиями к конфиденциальности: передать исходные данные третьей стороне
часто невозможно. Технологии многосторонних конфиденциальных вычислений (МКВ)
позволяют выполнять обучение и вывод над данными в виде долей секрета, не
раскрывая их отдельным участникам~\cite{zapechnikov2021two,zapechnikov2022three}.

В НИРС~1 были рассмотрены угрозы безопасности систем машинного
обучения~\cite{ilyushin2022attacks,papernot2018sok}. Настоящая работа
продолжает её: определяется, какие из этих угроз закрывают МКВ, и
проектируется система на их основе.

\textbf{Цель работы:} спроектировать защищённую систему машинного обучения с
применением МКВ и подготовить экспериментальную проверку.

\textbf{Задачи:}
\begin{enumerate}
    \item сопоставить угрозы из НИРС~1 с возможностями МКВ;
    \item построить детальную модель нарушителя и угроз;
    \item разработать архитектуру системы и требования к ней;
    \item обосновать выбор протокола ABY3 и платформы TF Encrypted, параметров
    и способа вычисления решающего дерева;
    \item разработать методику экспериментов;
    \item выполнить пробный запуск на одной машине.
\end{enumerate}

    \section{Угрозы и область применимости МКВ}
\label{sec:threats}

\subsection{Угрозы, выявленные в НИРС 1}
% ЧЕРНОВИК: перенести краткий перечень из НИРС 1.
Угрозы систем машинного обучения группируются по этапам конвейера: данные
обучения, модель, данные эксплуатации~\cite{ilyushin2022attacks,papernot2018sok}.

\subsection{Угрозы, которые закрывают МКВ}
\begin{itemize}
    \item раскрытие обучающих данных вычислительным серверам;
    \item раскрытие весов модели при обучении и выводе;
    \item раскрытие входных данных клиента при выводе;
    \item компрометация одного сервера из трёх (при честном большинстве).
\end{itemize}

\subsection{Угрозы, которые МКВ не закрывают}
\begin{itemize}
    \item атаки на результат: определение принадлежности к обучающей
    выборке~\cite{shokri2017mia}, обращение и извлечение модели;
    \item отравление данных владельцем данных, честно выполняющим протокол;
    \item состязательные примеры;
    \item сговор двух серверов из трёх;
    \item побочные каналы, отказ в обслуживании, ошибки реализации.
\end{itemize}

\subsection{Сводное сопоставление}
\begin{table}[H]
    \caption{Сопоставление угроз и возможностей МКВ (заготовка)}
    \label{tab:threat-map}
    \centering
    \begin{tabular}{|p{6cm}|p{3cm}|p{5.5cm}|}
        \hline
        Угроза & Закрывает МКВ & Дополнительная мера \\ \hline
        Раскрытие обучающих данных серверу & да & --- \\ \hline
        Определение принадлежности к выборке & нет & дифференциальная
        приватность, ограничение запросов \\ \hline
        Отравление данных & нет & валидация входных данных \\ \hline
        Сговор большинства серверов & нет & организационное разнесение
        серверов \\ \hline
    \end{tabular}
\end{table}

\subsection{Выводы по разделу}
% ЧЕРНОВИК

    \section{Модель нарушителя и угроз}
\label{sec:adversary}

\subsection{Участники системы и их роли}
% ЧЕРНОВИК
Владельцы данных, три вычислительных сервера, владелец модели, получатель
результата, администратор.

\subsection{Классификация нарушителей}
Внешний и внутренний; получестный и активный; статический и адаптивный.
Порог компрометации для ABY3~--- один сервер из трёх~\cite{mohassel2018aby3,araki2016}.

\subsection{Возможности и цели нарушителя}

\subsection{Допущения безопасности}
Защищённые каналы между серверами, независимое администрирование серверов,
корректная генерация случайности.

\subsection{Формальная постановка}
Безопасность определяется в парадигме <<реальный мир~--- идеальный мир>>.

\subsection{Перечень актуальных угроз}
% Оформить по Методике оценки угроз ФСТЭК России (2021).

\subsection{Остаточные риски}
Методика моделирования атак на модели машинного обучения~\cite{podtopelny2024}.

    \section{Архитектура и требования}
\label{sec:architecture}

\subsection{Функциональные требования}
% ЧЕРНОВИК
Загрузка данных в виде долей, обучение, вывод, выдача результата
авторизованной стороне.

\subsection{Требования безопасности}

\subsection{Нефункциональные требования}
Допустимая потеря точности, время обучения и вывода, сетевой трафик.

\subsection{Архитектура системы}
Обобщенная модель архитектуры конфиденциального машинного обучения
приведена в~\cite{zapechnikov2021two}.

% Пример вставки рисунка (файл положить в inc/):
% \begin{figure}[H]
%     \centering
%     \includegraphics[width=0.9\linewidth]{inc/architecture.pdf}
%     \caption{Архитектура системы}
%     \label{fig:architecture}
% \end{figure}

\subsection{Жизненный цикл данных}
Разделение $\rightarrow$ вычисление $\rightarrow$ восстановление.

\subsection{Трассировка требований}

    \section{Протокол и платформа}
\label{sec:protocol}

\subsection{Обзор альтернатив}
% ЧЕРНОВИК
SecureML~\cite{mohassel2017secureml}, SecureNN~\cite{wagh2019securenn},
Falcon~\cite{wagh2021falcon}, CrypTen~\cite{knott2021crypten}; обзоры
трёх- и четырёхсторонних систем~\cite{zapechnikov2022three,zapechnikov2022four}.
Большинство этих работ ориентировано на нейронные сети, где узким местом
являются умножения и аппроксимации~\cite{lapina2024}; для решающего дерева
узкое место~--- сравнения.

\subsection{Обоснование выбора протокола ABY3}
ABY3~\cite{mohassel2018aby3} использует реплицированное разделение секрета
на кольце $\mathbb{Z}_{2^{64}}$~\cite{araki2016} и переходы между
арифметическим, булевым и Yao-представлениями.

\subsection{Обоснование выбора платформы TF Encrypted}
TF Encrypted~\cite{dahl2018tfe} предоставляет Keras-подобный интерфейс и
реализацию ABY3.

\subsection{Параметры протокола}
Значение $x$ кодируется в фиксированной точке с $f$ дробными битами:
\begin{equation}
    \label{eq:fixed-point}
    \tilde{x} = \lfloor x \cdot 2^{f} \rceil \bmod 2^{64}.
\end{equation}

Для решающего дерева в фиксированной точке кодируются признаки объекта и
пороги в узлах; разрядность $f$ должна сохранять порядок значений, иначе
меняется результат сравнения.

\subsection{Вычисление решающего дерева в МКВ}
% ЧЕРНОВИК
Модель~--- решающее дерево классификации глубины $d$. В отличие от
нейронных сетей, аппроксимации нелинейных функций (сигмоиды, softmax) не
требуются: основная операция~--- конфиденциальное сравнение признака с
порогом в каждом узле. В ABY3 сравнение выполняется переходом от
арифметического представления к булевому и извлечением старшего
бита~\cite{mohassel2018aby3}.

Чтобы путь по дереву не раскрывал данные, вычисляются результаты
сравнений во всех $2^{d}-1$ внутренних узлах, после чего значение листа
выбирается конфиденциально (произведение битов сравнения вдоль пути).
Стоимость растёт экспоненциально по $d$, поэтому глубина дерева~---
ключевой параметр экспериментов.

В TF Encrypted нет готовой реализации решающего дерева: её предстоит
собрать из примитивов сравнения и умножения протокола ABY3. Это главный
технический риск, который проверяется при пробном запуске
(раздел~\ref{sec:trial}).

    \section{Методика экспериментов}
\label{sec:methodology}

\subsection{Объект сравнения и контрольные вычисления}

Цель эксперимента состоит в проверке двух свойств системы: корректности
конфиденциального вывода по заданному дереву и зависимости затрат на вывод
от размера дерева и пакета запросов. Единицей сравнения служит прогноз для
одного объекта. Во всех вариантах используются одна версия дерева, один
порядок признаков и одно правило перехода в левое поддерево при равенстве
признака и порога. Обучение завершается до начала измерений.

Для объекта $x_i$ регистрируются три метки класса. Через $y_i^{(0)}$
обозначается результат открытого обхода исходного дерева, через $y_i^{(q)}$
--- результат открытого обхода после кодирования признаков и порогов
по формуле~\eqref{eq:fixed-point}, через $y_i^{(s)}$~--- восстановленная
метка конфиденциального вычисления. В целочисленном эталоне применяется
обычный последовательный обход по выбранной ветви. Это позволяет проверить
вычисление всех узлов и скрытый выбор листа по независимой реализации
той же функции.

На наборе из $n$ объектов вычисляются две доли расхождений:
\[
  e_{\mathrm{impl}}=\frac{1}{n}\sum_{i=1}^{n}
       [y_i^{(s)}\ne y_i^{(q)}], \qquad
  e_{\mathrm{quant}}=\frac{1}{n}\sum_{i=1}^{n}
       [y_i^{(q)}\ne y_i^{(0)}].
\]
Квадратные скобки обозначают индикатор события. Критерий корректности
реализации на контрольном наборе имеет вид $e_{\mathrm{impl}}=0$.
Величина $e_{\mathrm{quant}}$ характеризует влияние представления чисел
на итоговый прогноз. При наличии истинных меток дополнительно определяется
точность классификации каждого варианта по отношению к разметке. Совпадение
двух реализаций и правильность прогноза относительно разметки отвечают
на разные экспериментальные вопросы.

\subsection{Формирование контрольных наборов}

Для проверки сравнений формируется синтетический набор с заранее
известными путями от корня до листьев. Для каждого проверяемого порога
выбираются значения по обе стороны от него, точное равенство и пара
различных исходных значений с одинаковым кодом. Включаются отрицательные
признаки и пороги. Для проверки внутреннего узла остальные признаки
подбираются так, чтобы открытый обход достигал этого узла. В тестовом
дереве листьям присваиваются различные коды: тогда ошибка выбора пути
проявляется в итоговой метке. Отдельный тест с повторяющимися классами
проверяет, что получателю выдаётся код класса, предусмотренный моделью.

Для исследования на размеченных данных предусматривается табличная
выборка с числовыми признаками. До обучения фиксируются версия набора,
правила обработки пропусков, разбиение на обучающую и контрольную части
и начальное состояние генератора, используемого при разбиении.
Преобразования признаков определяются по обучающей части и затем
применяются к контрольной. Для каждого обученного дерева сохраняются
топология, индексы признаков, пороги и метки листьев. Это обеспечивает
повторное использование одной модели во всех трёх вычислениях.

\subsection{План измерений}

Варьируются глубина дерева $d$, фактическое число внутренних узлов $m$,
размер пакета $B$ и число дробных разрядов $f$. Для оценки влияния
округления изменяют $f$ при фиксированных дереве и входах. Для оценки
производительности фиксируют $f$ и повторяют вывод на одном и том же
пакете с различными конфигурациями дерева. Наряду с глубиной указывается
$m$, поскольку деревья одинаковой глубины могут требовать разного числа
сравнений. До начала серии проверяется представимость значений и разностей,
используемых для знакового сравнения.

Время подготовки и загрузки модели учитывается отдельно. Задержка пакета
измеряется от начала формирования долей подготовленных входов до
восстановления последней метки у получателя. При длительности обработки
$t_B$ пропускная способность равна $B/t_B$. Первый запуск регистрируется
отдельно; после прогрева выполняются повторные измерения при фиксированных
версиях программных компонентов и одинаковом выделении ресурсов.
Для каждой конфигурации приводятся число повторений, медиана задержки
и межквартильный размах.

Для каждого вычислительного процесса регистрируются процессорное время
и пиковое потребление оперативной памяти. При наличии инструментированного
межсерверного транспорта объём обмена определяется суммой отправленных
байтов по трём серверам, без повторного учёта при приёме. Число раундов
определяется по последовательным этапам протокольного взаимодействия;
несколько сообщений одного этапа могут передаваться параллельно.
Если локальный режим исполняет операции в общей среде без наблюдаемого
сетевого обмена, сетевые показатели для такой конфигурации обозначаются
как недоступные. Переход к измерению производительности предусмотрен
после прохождения функциональных проверок раздела~\ref{sec:trial}.


Для первой серии выбираются полные бинарные деревья глубины
$d\in\{1,2,4,6\}$ и пакеты $B\in\{1,16,64\}$. Декартово произведение
этих множеств задаёт 12 конфигураций; значения $m$ равны соответственно
1, 3, 15 и 63. Во всех конфигурациях используется единое $f$,
прошедшее проверку диапазона. Для каждой конфигурации планируется один
прогревочный запуск и 20 измеряемых повторов на одном фиксированном пакете.
В отдельной серии на фиксированных дереве и наборе объектов изменяется
только $f$; выбор допустимых значений определяется интервалами признаков
и порогов.

    \section{Пробный запуск: постановка и порядок проверки}
\label{sec:trial}

\subsection{Назначение локального стенда}

Пробный запуск предназначен для проверки согласованности трёх компонентов:
кодирования входов, выполнения операций ABY3 в TF~Encrypted и
восстановления метки класса. Предусматривается размещение трёх логических
вычислительных сторон на одном компьютере. Такой стенд позволяет сверять
результат после каждого тестового сценария. Независимое администрирование
трёх серверов, предусмотренное моделью нарушителя, требует отдельной
распределённой конфигурации.

Для воспроизведения запуска в протокол испытания включаются сведения
об операционной системе, архитектуре процессора, объёме памяти, версиях
Python, TensorFlow и TF~Encrypted. Перед загрузкой дерева требуется
проверить доступность реализации ABY3, назначение трёх различных ролей
и согласованность типов входных и выходных значений. Масштабирование
фиксируется в одном месте: при передаче уже закодированного целого значения
повторное умножение на масштаб изменило бы операнд сравнения.

\subsection{Контрольные сценарии}

Сначала проверяются операции над разделёнными тестовыми значениями.
Для операндов $-2$ и $3$ ожидаемые результаты сложения и умножения равны
соответственно $1$ и $-6$. Далее проверяется узел с порогом $0$:
при входах $-1$, $0$ и $1$ условие перехода в левую ветвь должно дать
значения $1$, $1$ и $0$. Для каждого сценария восстанавливается только
оговорённый тестовый выход. Эти проверки локализуют ошибки знаковой
интерпретации и обработки равенства до исполнения полного дерева.

Следующий сценарий задаёт дерево глубины два с двумя числовыми
признаками. В корне проверяется условие $x_1\leq 0$. В левом поддереве
проверяется $x_2\leq -1$ с кодами листьев $0$ и $1$, в правом~---
$x_2\leq 1$ с кодами $2$ и $3$. При $f=2$ все указанные пороги и целые
тестовые входы представимы точно. Для пакета
\[
  ((-1,-2),\ (0,0),\ (1,1),\ (1,2))
\]
ожидаемая последовательность меток равна $(0,1,2,3)$. Последовательность
должна сохраняться при поочерёдной подаче тех же объектов. Повторная
обработка пакета со свежими долями должна давать те же открытые метки.

Отдельный пример проверяет влияние округления. Для входа
$(1/16,0)$ открытый обход исходного дерева даёт класс $2$. При $f=2$
первая координата округляется к нулю, поэтому целочисленный эталон
даёт класс $1$. Ожидаемый результат ABY3 здесь равен $1$: его совпадение
с целочисленным эталоном одновременно с расхождением с исходной моделью
проверяет трактовку двух показателей раздела~\ref{sec:methodology}.

\subsection{Регистрация результата}

Для каждого сценария предусматривается запись версии модели, значения
$f$, открытого синтетического входа, ожидаемой и фактически полученной
меток. При расхождении фиксируется первый сценарий, на котором оно
возникло; для диагностики отдельно повторяются сравнение и выбор
дочернего значения на тех же синтетических операндах. Критерием успешного
пробного запуска служит совпадение всех фактических выходов с указанными
эталонами, включая порядок меток в пакете. Приведённые числовые значения
получены аналитически и предназначены для последующей сверки при исполнении
прототипа.

    \conclusion

% ЧЕРНОВИК: перечислить результаты по каждой задаче из введения.
В ходе работы были получены следующие результаты.


    \printbibliography

    \appendix
    \appendixsection{Полная таблица угроз}
\label{app:threats}

\begin{table}[htbp]
\caption{Реестр сценариев для проектируемой системы}
\centering\small
\begin{tabularx}{\linewidth}{|p{.08\linewidth}|X|X|}
\hline
Код & Источник, условие и объект & Последствие \\ \hline
У1 & Совместный доступ к состояниям двух узлов & Восстановление признаков или параметров дерева \\ \hline
У2 & Наблюдение ветвления по адресу следующего узла & Раскрытие исхода сравнения \\ \hline
У3 & Подмена узла или смешение версий при передаче долей & Раскрытие доли либо ошибочный прогноз \\ \hline
У4 & Получение результата посторонней стороной или накопление разрешённых ответов & Раскрытие метки или сведений о модели \\ \hline
У5 & Потеря вычислительного узла во время запроса & Нарушение доступности \\ \hline
У6 & Недостоверная модель или искажённый вход & Нарушение целостности решения \\ \hline
\end{tabularx}
\end{table}

\appendixsection{Конфигурация окружения пробного запуска}
\label{app:env}

Дата и идентификатор сборки, ОС, процессор, память, версии Python, TensorFlow и TF~Encrypted, режим транспорта, тайм-ауты и расположение журнала заполняются по фактическому стенду.

\end{document}
